\section{2021: \work{Traditionis custodes} and the Problem of Article 8}

\subsection{The act}

On 16 July 2021 Francis issued the apostolic letter given motu proprio
\work{Traditionis custodes}, promulgated by publication in
\work{L'Osservatore Romano} and entering immediately into force. Four of
its eight articles bear directly on the Fraternity of Saint Peter, and
the fourth of those bears on it fatally unless something else
intervenes.\endnote{Francis, apostolic letter given motu proprio
\work{Traditionis custodes}, 16 July 2021, quoted from the Holy See's own
English version
(\url{https://www.vatican.va/content/francesco/en/motu_proprio/documents/20210716-motu-proprio-traditionis-custodes.html}),
read 2026-07-25; the exact bytes read on that date match the artifact
registered in this repository's source library. Latin in \work{Acta
Apostolicae Sedis} 113 (2021) 793--796 as cited in the canonical
literature; the 2021 \work{Acta} volume is not available in the Holy See's
whole-volume PDF archive and the locus was not independently checked.}

\emph{Art.~1}: ``The liturgical books promulgated by Saint Paul VI and
Saint John Paul II, in conformity with the decrees of Vatican Council II,
are the unique expression of the \latin{lex orandi} of the Roman Rite.''
The change of one adjective---\latin{unica} for \work{Summorum
Pontificum}'s ordinary-and-extraordinary---undoes the two-usages doctrine
on which the Fraternity's public self-description had rested since 2007.

\emph{Art.~2}: it belongs to the diocesan bishop to regulate the
liturgical celebrations of his diocese, and ``it is his exclusive
competence to authorize the use of the 1962 Roman Missal in his diocese,
according to the guidelines of the Apostolic See.''

\emph{Art.~6}: ``Institutes of consecrated life and Societies of
apostolic life, erected by the Pontifical Commission \work{Ecclesia Dei},
fall under the competence of the Congregation for Institutes of
Consecrated Life and Societies for Apostolic Life.''

\emph{Art.~8}: ``Previous norms, instructions, permissions, and customs
that do not conform to the provisions of the present Motu Proprio are
abrogated.''

\subsection{Two very different articles}

Article 6 is the third act in the sequence that began in 1988. Faculty 6
of \latin{Quia peculiare munus} had reserved the exercise of the Holy
See's authority over the erected societies to the Commission
\latin{donec aliter provideatur}; the motu proprio of 2019 moved that
exercise to the Congregation for the Doctrine of the Faith; and
\work{Traditionis custodes} art.~6 moves the institutes and societies
themselves to the dicastery for consecrated life. Two years and six months
separate the two transfers. Neither act mentions any institute by name;
both act on the class defined by the erecting organ.

Article 8 is a different kind of provision entirely, and this account
does not pretend that its application to the Fraternity is obvious. Read
at its widest, it abrogates the concession of the 1962 books made by the
decree of 10 September 1988 and repeated in the decree of erection, since
those are ``previous permissions'' and they do not conform to a motu
proprio whose art.~2 gives the diocesan bishop exclusive competence to
authorize the 1962 Missal. Read more narrowly, an erection decree is a
constitutive act rather than a norm, instruction, permission or custom,
and what it constitutes---an institute whose stated object includes the
observance of certain traditions---is not the kind of thing art.~8
abrogates. The Fraternity's own communiqué of 20 July 2021 does not argue
either way; it says something else.

\subsection{The Fraternity's answer, 20 July 2021}

Four days after the motu proprio the Fraternity issued an official
communiqué from Fribourg. Its argument is not canonical but
characterological, and its tone is worth recording exactly, because it is
the only public statement the institute has ever made about a papal act
that touched its reason for existing.

It says the Fraternity ``has received Pope Francis' Motu Proprio
\work{Traditionis Custodes} with surprise.'' It recalls that the
Fraternity was ``founded and canonically approved according to the Motu
Proprio \work{Ecclesia Dei Adflicta}'' and ``has always professed its
adherence to the entire Magisterium of the Church and its fidelity to the
Roman Pontiff and the successors of the Apostles, exercising its ministry
under the responsibility of the diocesan bishops.'' It states that the
Fraternity ``is deeply saddened by the reasons given for limiting the use
of the Missal of Pope St.\ John XXIII, which is at the center of its
charism,'' and---the sentence that carries the weight---``The Fraternity in
no way recognizes itself in the criticisms made.'' It asks how one can
fail to notice the fruits of the apostolates attached to that Missal and
the youth of the communities attached to it. And it closes by reaffirming
``our unwavering fidelity to the successor of Peter on the one hand, and
on the other, our desire to remain faithful to our Constitutions and
charism, continuing to serve the faithful as we have done since our
foundation,'' hoping ``to be able to count on the understanding of the
bishops, whose authority we have always respected.''\endnote{``Official
communiqué following the publication of Traditionis Custodes''
(fssp.org), Fribourg, 20 July 2021, read 2026-07-25. Party statement;
quotations are from the Fraternity's own English text. Note its citation
of its own constitutions n.~8 and of Benedict XVI's address to the Roman
Curia of 22 December 2005 on ``the hermeneutic of reform in the continuity
of the Church.''}

Two features are worth drawing out. The communiqué locates the
Fraternity's defence in its constitutions---``our Constitutions and
charism''---and not in the 1988 liturgical concession; that is the choice
of the stronger ground, since a constitution definitively approved by the
Apostolic See is protected by the just autonomy of canon 586 as applied by
canon 732, whereas a permission is exactly what art.~8 abrogates. And it
announces no change of practice, asks for no clarification in terms, and
lodges no recourse. It waits.

\subsection{The \work{Responsa ad dubia} of 4 December 2021}

On 4 December 2021 the Congregation for Divine Worship and the Discipline
of the Sacraments published responses to eleven questions on the
application of \work{Traditionis custodes}, addressed to the presidents of
the conferences of bishops, with the pope's assent given in the audience
of 18 November 2021. Two of them matter here.

The first concerns places: a diocesan bishop may ask the Congregation for
a dispensation from art.~3 §\,2 to allow celebration in a parish church,
but only if it is established that no other church, oratory or chapel can
be used, ``the assessment of this impossibility'' being made ``with the
utmost care''; the celebration is not to appear in the parish Mass
schedule and is to be withdrawn when another venue becomes
available.\endnote{Congregation for Divine Worship and the Discipline of
the Sacraments, \work{Responsa ad dubia} on certain provisions of
\work{Traditionis custodes}, 4 December 2021, quoted from the Holy See's
own English text
(\url{https://www.vatican.va/roman_curia/congregations/ccdds/documents/rc_con_ccdds_doc_20211204_responsa-ad-dubia-tradizionis-custodes_en.html}),
read 2026-07-25; the exact bytes read match the artifact registered in
this repository's source library.}

The second concerns books: asked whether the sacraments may be celebrated
with the pre-conciliar \work{Rituale Romanum} and \work{Pontificale
Romanum}, the answer is ``Negative,'' and the explanation is that ``the
diocesan Bishop is authorised to grant permission to use only the
\work{Rituale Romanum} (last \latin{editio typica} 1952) and not the
\work{Pontificale Romanum} which predate the liturgical reform,'' and may
grant even that only to canonically erected personal parishes celebrating
with the 1962 Missal.

Set that beside the decree of 10 September 1988, which conceded to the
Fraternity the Missal, the Ritual, the \emph{Pontifical} and the
Breviary, and the shape of the coming difficulty is plain. The
\work{Responsa} do not name the Fraternity, and they are framed
throughout as answers about what a \emph{diocesan bishop} may permit. They
therefore do not, on their face, decide anything about a faculty granted
to a society of apostolic life by the Apostolic See. But they establish
what the general law now allows a bishop to allow, and a great deal of the
Fraternity's practical life happens in dioceses.

\subsection{Art.~6 in operation}

One consequence of art.~6 was immediate and administrative and is
recorded by the Fraternity itself. Writing in September 2024 it observed
that ``the Dicastery for Institutes of Consecrated Life and Societies of
Apostolic Life has been in charge of the FSSP and other former Ecclesia
Dei institutes for the past three years,'' dating its own change of
supervisor to 2021 and to no other cause.\endnote{``Apostolic visitation
of the Fraternity'' (fssp.org), 26 September 2024, read 2026-07-25. The
Fraternity's arithmetic---``the past three years''---places the transfer in
2021, that is, at \work{Traditionis custodes} art.~6 and not at the
decommission of 2019, which had moved the Commission's tasks to the
doctrinal congregation.}

\actrecord{\work{Traditionis custodes}, 16 July 2021, arts.~1, 2, 6 and
8; \work{Responsa ad dubia}, 4 December 2021.}{That the reformed books
became the unique expression of the \latin{lex orandi} of the Roman Rite;
that authorizing the 1962 Missal in a diocese became the exclusive
competence of its bishop under Apostolic See guidelines; that institutes
and societies erected by the Commission \work{Ecclesia Dei} passed to the
dicastery for consecrated life; that previous non-conforming norms,
instructions, permissions and customs were abrogated; and that a diocesan
bishop may permit the 1952 Ritual only to personal parishes and may not
permit the pre-conciliar Pontifical at all.}{Neither act names the
Fraternity of Saint Peter or any other institute. The \work{Responsa}
concern the powers of diocesan bishops. Whether art.~8 reaches a faculty
granted directly by the Apostolic See to a society of apostolic life is a
question these texts do not answer, and this account does not answer it
either.}
